\section{Results}
\label{sec:results}

\subsection{Cross-well low-data results}

Table~\ref{tab:crosswell_clp_vs_ae} compares CLP-CSGM Ridge against the
strongest direct baseline, AE \([u,b]\), for every low-data and sparse
measurement cell. CLP-CSGM Ridge wins in all nine cells. The relative RMSE gap
is modest at step 8 and \(\rho=0.05\), but becomes large as the low-data regime
becomes more severe, reaching improvements above 38\% for step 32 at
\(\rho=0.10\) and \(\rho=0.20\).

\begin{table}[htbp]
  \centering
  \caption{Cross-well Vc comparison between CLP-CSGM Ridge and AE \([u,b]\).
  Negative gap means CLP-CSGM has lower RMSE.}
  \label{tab:crosswell_clp_vs_ae}
  \small
  \begin{tabular}{rrrrrr}
\toprule
step & n\_train\_approx & rho & clp\_csgm\_rmse & ae\_ub\_rmse & gap\_vs\_ae\_pct \\
\midrule
8 & 104 & 0.05000 & 0.01541 & 0.01709 & -9.84267 \\
8 & 104 & 0.10000 & 0.01432 & 0.01635 & -12.40651 \\
8 & 104 & 0.20000 & 0.01167 & 0.01633 & -28.52507 \\
16 & 52 & 0.05000 & 0.01493 & 0.01891 & -21.01886 \\
16 & 52 & 0.10000 & 0.01391 & 0.01803 & -22.81414 \\
16 & 52 & 0.20000 & 0.01247 & 0.01808 & -31.00344 \\
32 & 27 & 0.05000 & 0.01562 & 0.02132 & -26.71650 \\
32 & 27 & 0.10000 & 0.01484 & 0.02501 & -40.67353 \\
32 & 27 & 0.20000 & 0.01259 & 0.02035 & -38.15454 \\
\bottomrule
\end{tabular}

\end{table}

\begin{figure}[htbp]
  \centering
  \includegraphics[width=\linewidth]{results_crosswell_clp_vs_ae.png}
  \caption{Cross-well Vc RMSE for CLP-CSGM Ridge and AE \([u,b]\) across the
  low-data grid. Each panel fixes the sliding-window step and varies the sparse
  measurement ratio.}
  \label{fig:crosswell_clp_vs_ae}
\end{figure}

Figure~\ref{fig:crosswell_all_methods} places this pairwise comparison in the
full benchmark context. The direct MLP and PCA baselines become unstable in the
low-data cells, while ML only remains more stable but less accurate. The main
competitor is AE \([u,b]\), and CLP-CSGM Ridge consistently stays below it
across the measurement ratios. This supports the claim that the gain is not an
artifact of comparing against weak baselines; it is visible against the
strongest direct representation-learning baseline.

\begin{figure}[htbp]
  \centering
  \includegraphics[width=\linewidth]{results_crosswell_rmse_vs_rho.png}
  \caption{Cross-well Vc RMSE for all main benchmarks. CLP-CSGM Ridge remains
  the best method across the low-data grid, while direct MLP and PCA baselines
  degrade sharply when training windows are scarce.}
  \label{fig:crosswell_all_methods}
\end{figure}

The overall ranking in Table~\ref{tab:crosswell_rank} shows the same pattern.
Across the nine step--\(\rho\) cells, CLP-CSGM Ridge obtains mean RMSE
0.01397, compared with 0.01905 for AE \([u,b]\). This corresponds to a relative
gap of about 26.65\% against the strongest direct baseline.

\begin{table}[htbp]
  \centering
  \caption{Overall cross-well Vc ranking across all low-data cells.}
  \label{tab:crosswell_rank}
  \small
  \begin{tabular}{lrrr}
\toprule
method & mean\_rmse & mean\_relative\_l2 & cells \\
\midrule
CLP-CSGM Ridge & 0.01397 & 0.18131 & 9 \\
AE [u,b] & 0.01905 & 0.25464 & 9 \\
ML only & 0.02403 & 0.32932 & 9 \\
PCA [u,b] & 0.14654 & 2.08220 & 9 \\
MLP [u,b] & 0.20112 & 2.79301 & 9 \\
\bottomrule
\end{tabular}

\end{table}

\subsection{Real-well F03 GR-only results}

Table~\ref{tab:f03_results} reports the F03-4 GR-only porosity benchmark.
CLP-CSGM Ridge has the best mean RMSE across measurement ratios, 0.01360 versus
0.01398 for AE \([u,b]\). The gain is concentrated in the sparse-measurement
regime: CLP-CSGM Ridge is clearly better at \(\rho=0.20\) and \(\rho=0.30\).
At higher measurement ratios, AE \([u,b]\) becomes slightly better.

\begin{table}[htbp]
  \centering
  \caption{Real-well F03-4 GR-only porosity benchmark by measurement ratio.}
  \label{tab:f03_results}
  \small
  \begin{tabular}{rlrrr}
\toprule
rho & method & rmse & mae & relative\_l2 \\
\midrule
0.20000 & CLP-CSGM Ridge & 0.01575 & 0.01365 & 0.05909 \\
0.20000 & AE [u,b] & 0.01797 & 0.01560 & 0.06737 \\
0.20000 & PCA [u,b] & 0.02079 & 0.01712 & 0.07791 \\
0.20000 & ML only & 0.02546 & 0.02219 & 0.09550 \\
0.20000 & MLP [u,b] & 0.03189 & 0.02731 & 0.11949 \\
0.30000 & CLP-CSGM Ridge & 0.01449 & 0.01245 & 0.05437 \\
0.30000 & AE [u,b] & 0.01536 & 0.01330 & 0.05760 \\
0.30000 & PCA [u,b] & 0.02241 & 0.01834 & 0.08398 \\
0.30000 & ML only & 0.02546 & 0.02219 & 0.09550 \\
0.30000 & MLP [u,b] & 0.03184 & 0.02716 & 0.11924 \\
0.40000 & AE [u,b] & 0.01296 & 0.01107 & 0.04856 \\
0.40000 & CLP-CSGM Ridge & 0.01342 & 0.01141 & 0.05034 \\
0.40000 & PCA [u,b] & 0.02206 & 0.01790 & 0.08263 \\
0.40000 & ML only & 0.02546 & 0.02219 & 0.09550 \\
0.40000 & MLP [u,b] & 0.03728 & 0.03245 & 0.13966 \\
0.50000 & AE [u,b] & 0.01234 & 0.01065 & 0.04625 \\
0.50000 & CLP-CSGM Ridge & 0.01246 & 0.01061 & 0.04673 \\
0.50000 & PCA [u,b] & 0.02358 & 0.01909 & 0.08834 \\
0.50000 & ML only & 0.02546 & 0.02219 & 0.09550 \\
0.50000 & MLP [u,b] & 0.04466 & 0.03965 & 0.16717 \\
0.60000 & AE [u,b] & 0.01126 & 0.00955 & 0.04218 \\
0.60000 & CLP-CSGM Ridge & 0.01186 & 0.01006 & 0.04447 \\
0.60000 & PCA [u,b] & 0.02417 & 0.01960 & 0.09056 \\
0.60000 & ML only & 0.02546 & 0.02219 & 0.09550 \\
0.60000 & MLP [u,b] & 0.04423 & 0.03873 & 0.16568 \\
\bottomrule
\end{tabular}

\end{table}

\begin{figure}[htbp]
  \centering
  \includegraphics[width=0.92\linewidth]{01_rmse_vs_measurement_ratio.png}
  \caption{Real-well F03-4 GR-only porosity RMSE by measurement ratio. CLP-CSGM
  Ridge is strongest at the sparsest ratios, while AE \([u,b]\) becomes
  competitive as more target measurements are available.}
  \label{fig:f03_gr_only_rmse}
\end{figure}

Figure~\ref{fig:f03_all_methods} shows the same F03 result with all main
benchmarks. CLP-CSGM Ridge dominates the sparse regime, while AE \([u,b]\)
overtakes it at higher measurement ratios. The crossover is informative rather
than problematic: it marks the regime where enough target-well samples are
available for a direct feature-based model to be competitive.

\begin{figure}[htbp]
  \centering
  \includegraphics[width=0.92\linewidth]{05_rmse_grouped_bars_by_ratio.png}
  \caption{Real-well F03-4 GR-only RMSE for all main benchmarks. CLP-CSGM Ridge
  has its strongest advantage at low measurement ratios, whereas AE \([u,b]\)
  becomes slightly better when more sparse target measurements are available.}
  \label{fig:f03_all_methods}
\end{figure}

The MAE and relative-\(\ell_2\) curves in
Figs.~\ref{fig:f03_mae}--\ref{fig:f03_rell2} confirm that the RMSE behavior is
not metric-specific. CLP-CSGM Ridge stays close to or below AE \([u,b]\) across
the sparse-to-moderate ratios, while the direct MLP baseline deteriorates as
more sparse-measurement features are added. This suggests that naively
concatenating \(b\) to \(u\) can inject high-variance features, whereas the
proposed method uses \(b\) only through the physical measurement residual.

\begin{figure}[htbp]
  \centering
  \includegraphics[width=0.92\linewidth]{02_mae_vs_measurement_ratio.png}
  \caption{Real-well F03-4 GR-only MAE by measurement ratio. The ranking is
  consistent with RMSE: CLP-CSGM Ridge is strongest in the sparse regime and
  remains competitive as the measurement ratio increases.}
  \label{fig:f03_mae}
\end{figure}

\begin{figure}[htbp]
  \centering
  \includegraphics[width=0.92\linewidth]{03_relative_l2_vs_measurement_ratio.png}
  \caption{Relative \(\ell_2\) error on the real-well F03-4 GR-only benchmark.
  The relative-error view preserves the same conclusion as RMSE and MAE,
  indicating that the gain is not caused by scale effects alone.}
  \label{fig:f03_rell2}
\end{figure}

Figures~\ref{fig:f03_gain_rmse}--\ref{fig:f03_gain_mae} express the same
results as gain over the ML-only baseline. Positive values indicate improvement
over using GR alone. CLP-CSGM Ridge and AE \([u,b]\) are the only methods with
large positive gains across all ratios, but CLP-CSGM Ridge has the stronger gain
when measurements are sparse. The direct MLP \([u,b]\) curve is negative,
reinforcing that sparse measurements should be incorporated through a
measurement model rather than only as raw input features.

\begin{figure}[htbp]
  \centering
  \includegraphics[width=0.92\linewidth]{06_gain_rmse_over_ml_only.png}
  \caption{RMSE gain over the ML-only baseline on F03-4. CLP-CSGM Ridge gives
  the largest gain at low measurement ratios and remains consistently positive
  across the full range.}
  \label{fig:f03_gain_rmse}
\end{figure}

\begin{figure}[htbp]
  \centering
  \includegraphics[width=0.92\linewidth]{07_gain_mae_over_ml_only.png}
  \caption{MAE gain over the ML-only baseline on F03-4. The gain pattern mirrors
  the RMSE gain, supporting the robustness of the comparison across metrics.}
  \label{fig:f03_gain_mae}
\end{figure}

The qualitative diagnostics in Figs.~\ref{fig:f03_examples}--\ref{fig:f03_depth}
show how these aggregate gains appear in reconstructed profiles. The example
windows show that MLP \([u,b]\) can produce oscillatory outputs, while PCA and
AE are smoother. CLP-CSGM Ridge is also smooth but reacts to sparse constraints
through the latent optimization. The parity plot shows the highest pooled
correlation among the displayed models, and the depth profile indicates that
CLP-CSGM tracks the observed porosity trend without the high-frequency
instability of the direct MLP baseline.

\begin{figure}[htbp]
  \centering
  \includegraphics[width=\linewidth]{08_example_ground_truth_vs_models.png}
  \caption{Example F03-4 test windows. CLP-CSGM Ridge produces smooth
  reconstructions that follow the ground-truth trend, while direct MLP
  \([u,b]\) is visibly more oscillatory.}
  \label{fig:f03_examples}
\end{figure}

\begin{figure}[htbp]
  \centering
  \includegraphics[width=0.82\linewidth]{09_parity_ground_truth_vs_prediction.png}
  \caption{Pooled parity plots for F03-4 test windows. CLP-CSGM Ridge has the
  highest displayed correlation, indicating better ordering of high and low
  porosity values across the test block.}
  \label{fig:f03_parity}
\end{figure}

\begin{figure}[htbp]
  \centering
  \includegraphics[width=\linewidth]{10_depth_profile_porosity.png}
  \caption{Depth-profile reconstruction for the F03-4 test block at
  \(\rho=0.20\). CLP-CSGM Ridge follows the observed porosity trend while
  avoiding the unstable oscillations of direct MLP \([u,b]\).}
  \label{fig:f03_depth}
\end{figure}

This behavior is consistent with the design of CLP-CSGM. When measurements are
very sparse, the log-conditioned latent prior prevents the inverse problem from
fitting arbitrary sparse observations, while the measurement-consistency term
corrects the weak GR-only prior. When measurements become denser, a direct AE
baseline can use the \([u,b]\) feature vector effectively and the extra latent
optimization may offer less advantage.
